\begin{afterword}
\summaryhead

The post takes its title from Sun Tzu's advice to leave a surrounded enemy a way out. At a rationalists' gathering the author met a woman who rejected cryonics because the soul would not stay with the frozen body. The author advised her to picture the world without souls, and what she would do in it, before weighing the evidence for souls.

The principle is to leave yourself a line of retreat. A feared layoff is less frightening once you have counted your savings and checked the job market, and only then can you judge the risk fairly. Imagining a feared possibility as a thought experiment takes less courage than admitting it is likely, and the hope is that it makes the second step easier. Visualizing a belief is not conceding it.

The technique needs self-honesty: to name which ideas scare you, and to accept the premise for real. The author uses it personally. Theists who say morality is impossible without God would find, on honest reflection, that they would not start killing, and quadriplegics are less sad than they expect. What is true is already so; imagining it does no harm.

\responsehead

The advice is sensible and old, and the post gives little evidence that its new use works.

Its core is a remedy for fear. Seneca told a friend facing a lawsuit to assume the feared outcome would happen and to measure it in his own mind. The post's own example, the layoff, is the same exercise. What the post adds is the purpose: to lower the fear so that the probability can be judged fairly. That is a real addition. The post names no precedent.

Whether the addition works is not shown. The mechanism is offered as ``the hope.'' The support is the author's own use of the technique, with no case described, and one conversation whose outcome the post does not report. It ends by calling the technique ``powerful.''

The first illustration answers people who say morality is impossible without God by showing that they would not start ``slaughtering babies.'' That meets the claim as a forecast of behaviour, the form it took in the conversation. Philosophers usually make a different claim with the same words, about whether anything would be objectively right or wrong without God, and some atheists, Nietzsche among them, have accepted a version of it. The illustration does not bear on that claim, and the post does not say it does.

The second illustration, that quadriplegics are less sad than they expect, comes without a source. Research supports its direction: a review reports that patients often rate their quality of life higher than the public rates the same conditions. It also finds that adaptation is partial. In two national panels, a long-term disability brought lasting drops in happiness with little recovery. ``People do get over things'' holds only in part.

The restatement of Gendlin, that owning up to the truth does not make it worse, carries the problem noted under ``You Can Face Reality.'' The post's narrower claim, that merely picturing a world does no harm, does not depend on it.

In short: an old remedy for fear, turned to the new use of weighing evidence fairly, recommended on the author's own practice, and illustrated with two examples: one meets only the behavioural reading of the claim about God and morality, and research supports the other only in part.
\end{afterword}
